% Design record, copied verbatim from the transition stage's working directory (transition/dissertation_refs.tex): the dissertation passages the transition engine implements.

% Chapter 6 (monitored_deviations_chapter.tex), the seed-dependent system: forward equations for an impulse with the state displaced and the noise-states informed (privy) or not (naive). Lines 212-330.
\section{The seed-dependent system}
\label{sec:fb_system}

Fix a deviating player $i$ and a seed time $t$, with privy set $\mathcal P_i$ and naive set $\mathcal N_i$ as in \eqref{eq:PiNi_def}.  The state dynamics are \eqref{eq:baseline_dynamics}, with player $m$'s control entering through $B^X_m$.  The forward equations come first, then how a privy player responds to the seed, the backward system in operator form and in coefficient blocks, and last the case in which the deviating player responds to its own seed.

\subsection{Forward deviation equations}

The initial condition for a state impulse by player $i$ at time $t$ is
\begin{equation}\label{eq:seed_initial}
    \delta_i X_t=v,
    \qquad
    \delta_i w^k_t(u)=0\quad(k\in \mathcal N_i,\ 0\le u\le t).
\end{equation}
For $s\ge t$, the displaced state is
\begin{equation}
    \left(\delta_i X_s,\ \{\delta_i w^k_s(u):k\in \mathcal N_i,\ 0\le u\le s\}\right).
\end{equation}
The state equation is
\begin{equation}\label{eq:full_forward_state}
\boxed{
\frac{d}{ds}\delta_i X_s
=B^X_X(s)\delta_i X_s
+\sum_{k\in \mathcal N_i}\int_0^s B^X_k(s)\,D^k_s(u)\,\delta_i w^k_s(u)\,du
+\sum_{j\in \mathcal P_i}B^X_j(s)\,\delta_i D^j_s .
}
\end{equation}
For each naive player $k\in \mathcal N_i$, define the naive error variation
\begin{equation}\label{eq:error_variation_full}
    \delta_i\widetilde X^k_s
    :=\delta_i X_s-\int_0^s X(s,r)\,\delta_i w^k_s(r)\,dr .
\end{equation}
The kernel $X(s,r)$ is the impulse response of the state at $s$ to a primitive shock at $r$.  It carries no player index, since the dynamics are common; integrating it against the truth gives the state, and against a player's estimates gives that player's estimate of the state.  Player $k$'s unresolved-state kernel $\widetilde X^k_s(u)$ of Chapter~\ref{ch:forecasting} is the part of that response player $k$ has not yet resolved, and its innovation kernel is
\begin{equation}\label{eq:CY_def_full}
    \widetilde C^{k,Y}_s(u):=B^{Y,k}_X(s)\widetilde X^k_s(u).
\end{equation}
The naive filter variation has an interior revision equation
\begin{equation}\label{eq:full_forward_filter_revision}
\boxed{
    \partial_s\delta_i w^k_s(u)
    =\big(\widetilde C^{k,Y}_s(u)\big)^\top B^{Y,k}_X(s)\,\delta_i\widetilde X^k_s,
    \qquad 0\le u<s,
}
\end{equation}
and a birth term
\begin{equation}\label{eq:full_forward_filter_birth}
\boxed{
    \delta_i w^k_s(s)
    =E^{k\top}B^{Y,k}_X(s)\,\delta_i\widetilde X^k_s .
}
\end{equation}
Substituting \eqref{eq:error_variation_full} into \eqref{eq:full_forward_filter_revision}, the coefficient of $\delta_iw^k_s(r)$ in $\partial_s\delta_iw^k_s(u)$ is
\[
-\widetilde C^{k,Y}_s(u)^\top B^{Y,k}_X(s)X(s,r).
\]
This is the belief--belief block of the generator, and it is not symmetric in $(u,r)$: the lag $u$ being revised enters through the error kernel, the lag $r$ of the shock being estimated through the state kernel.  That asymmetry makes the generator non-self-adjoint.

\subsection{A privy player's response}

Fix a privy player $j\in \mathcal P_i$.  To compute player $j$'s response to the origin-$i$ seed, one tests whether $j$ would profitably deviate from its proposed response.  This introduces a second, hypothetical, $j$-origin spike.  The $i$-seed and this test $j$-seed are distinct labels.  In the gain defined below, the rows are indexed by $\mathcal N_j$, the players naive to $j$, because they come from player $j$'s own deviation problem; the columns by $\mathcal N_i$, the players naive to $i$, because they come from the already-realized $i$-origin displacement.

Let the closed-loop coefficients for an $a$-origin deviation be
\begin{equation}\label{eq:effective_coefficients_general}
\begin{aligned}
    A^{XX}_a(s)
    &:=B^X_X(s)+\sum_{q\in \mathcal P_a}B^X_q(s)\,D^{q\leftarrow a,X}_s,\\
    A^{XW^p}_a(s,u)
    &:=B^X_p(s)\,D^p_s(u)+\sum_{q\in \mathcal P_a}B^X_q(s)\,D^{q\leftarrow a,W^p}_s(u),
    \qquad p\in \mathcal N_a.
\end{aligned}
\end{equation}
These coefficients use the response kernels appropriate to the origin label $a$.  When $a=i$ they propagate the observed $i$-seed; when $a=j$ they are the row-side coefficients that price a hypothetical deviation by $j$.

The $i$-variation of player $j$'s closed-loop adjoint has rows $\delta_i H^{X,j}_s$ and $\delta_i H^{W^{k'},j}_s(u)$ for $k'\in \mathcal N_j$, dual to the row state $S^j$.  The next two subsections write it as one rectangular gain, first in operator form and then in coefficient blocks, with source $G_{\mathrm{eff}}^{j\leftarrow i}$.

Player $j$'s first-order condition, conditioned on its own information $\mathcal F^j_s$, gives its response.  Player $j$'s strategies are measurable in the factor $\mathcal Z_i$ (Remark~\ref{rem:zero_variance}), so the origin-$i$ seed is a known input to $j$.  Subtracting the condition player $j$ satisfies on the equilibrium path from the one it satisfies after the $i$-seed gives
\begin{equation}\label{eq:privy_projected_response_foc}
    0
    =G^{DD,j}(s)\,\delta_iD^j_s
     +G^{DX,j}(s)\,\delta_iX_s
     +\E\!\left[
        (B^X_j(s))^\top\delta_iH^{X,j}_s
        \mid \mathcal F^{j}_s
      \right].
\end{equation}
Player $j$ observes the $i$-seed, and the forward and adjoint systems are linear with deterministic coefficients, so the conditional expectation collapses:
\begin{equation}\label{eq:full_privy_response}
\boxed{
    \delta_iD^j_s=-\big(G^{DD,j}(s)\big)^{-1}\big[G^{DX,j}(s)\,\delta_iX_s+(B^X_j(s))^\top\delta_iH^{X,j}_s\big] .
}
\end{equation}

\begin{remark}[Two labels, no cross term]\label{rem:fixed_kernel}
The forward system is affine in the seeds, so the $i$- and $j$-labeled forward effects add, and the test uses the same gains as the ordinary $j$-origin deviation problem.  The $i$-seed shifts the value of player $j$'s adjoint while leaving that gain system intact.
\end{remark}

\subsection{Operator form and source}
\label{sec:filtered_operator}


The operator notation compresses the seed-dependent physical and noise-state equations above.  Stacked symbols collect Chapter~\ref{ch:forecasting} coordinates, and the rectangular gain $H^{j\leftarrow i}$ collects their adjoint responses into one backward solve.

An $i$-origin seed moves the noise-states of $\mathcal N_i$, while a hypothetical deviation by $j$ moves those of $\mathcal N_j$.  There is then no single displaced state, so the origin label indexes it.  For any origin label $a$, define the state attached to that origin by
\begin{equation}\label{eq:filtered_state_Sa}
    S^a_s
    :=
    \left(
        X_s,\,
        \{\widehat W^p_s(u):p\in \mathcal N_a,\ 0\le u<s\}
    \right).
\end{equation}
The superscript $a$ is the origin label, and the noise-state fields are those of the players naive to origin $a$.

For an $a$-origin deviation, write
\begin{equation}\label{eq:delta_Sa_def}
    \delta_a S^a_s
    :=
    \left(
        \delta_a X_s,\,
        \{\delta_a w^p_s(u):p\in \mathcal N_a\}

% Chapter 1 appendices (finite_horizon_appendices.tex): the forward state-filter system, lines 60-100 and 330-360.
responds at first order through~$\delta F^i$, changing which part of the fixed policy kernel the player's own information resolves.

\looseness=-1 The variation of the unresolved adjoint decomposes as
\begin{equation}\label{eq:delta_Htilde_channels}
\delta\widetilde H^{X,i}_t(s)
=
\underbrace{
-\int_s^t H^{X,i}_t(u)\,\delta F^i_t(u,s)\,du
\vphantom{\int_{t_0}^T}
}_{\text{(a) filtering channel}}
\;+\;
\underbrace{
\int_{t_0}^T\Phi^{HD,i}(t,\tau)\,
\widetilde{\delta D}^i_{W,t}(\tau,s)\,d\tau
}_{\text{(b) equilibrium channel}},
\end{equation}
where $\Phi^{HD,i}$ is the deterministic kernel through which
a control variation at time $\tau$ enters the time-$t$ adjoint,
and $\widetilde{\delta D}^i_W$ is the unresolved
component of the control variation induced by $\delta F^i$.

Through channel~(a) the player resolves more of the fixed adjoint from the perturbed signal, the standard value-of-information effect present in single-agent problems.
Channel~(b) arises because the same $\delta F^i$ changes the
projection, leaving part of $D^i_{W,\tau}$
unresolved by the player's current information.
The equilibrium channel operates through the information wedge and vanishes when
$\mathcal{V}^i_t = 0$ (exogenous signals,
Corollary~\ref{cor:exogenous_signal}).
Multi-player rational inattention couples the
filtering and control problems that separate in single-agent
settings: the marginal value of precision depends on the
equilibrium through the wedge.
Optimizing attention requires solving for the equilibrium first.

%%%%%%%%%%%%%%%%%%%%%%%%%%%%%%%%%%%%%%%%%%%%%%%%%%%%%%%%%%%%%%%%%%%%%%%%%%%%%%%

\paragraph{Local profile-space notation.} In the remainder of this appendix, bold symbols denote stacked policy profiles, and calligraphic and sans-serif symbols denote maps on the profile space.  These are local functional-analytic abbreviations and do not carry into the substantive chapters.

\section{Well-posedness of the deterministic impulse-response fixed point}
\label{sec:wellposedness}

\mathcal S_T:=\prod_{i=1}^n\mathcal S_T^i
\]
be such a realization, using a profile-independent observation or reference-
innovation coordinate system.  If profile-dependent noise-state coordinates
are used computationally, the coordinate-change terms are included when the
best-response derivative is formed.  Proposition~\ref{prop:global_best_response}
defines the exact best-response map
\[
\mathcal B(\mathbf s)
:=\bigl(\mathcal B_i(\mathbf s^{-i})\bigr)_{i=1}^n,
\qquad \mathbf s\in\mathcal S_T,
\]
and the equilibrium residual
\begin{equation}
\label{eq:exact_br_residual}
\mathcal R(\mathbf s):=\mathbf s-\mathcal B(\mathbf s).
\end{equation}
A zero of $\mathcal R$ is a fixed point of the exact linear-strategy
best-response map, and so a Nash equilibrium (Corollary~\ref{cor:nsl-br}).

The forward estimates below follow by Gr\"onwall arguments on the
forward system~\eqref{eq:Xbar_dyn}--\eqref{eq:F_dyn}, using the a priori
bounds of Lemmas~\ref{lem:F_bounded}--\ref{lem:Xtilde_dissipation} to
control the quadratic nonlinearity in~\eqref{eq:F_dyn}.
The adjoint estimates follow by backward Gr\"onwall on the linear adjoint
system
\eqref{eq:HX_process_backward_revised}--\eqref{eq:Hk_process_backward_revised},
using Cauchy--Schwarz on the belief-price coupling terms, including the birth term $B^{Y,k}_X(t)^\top E^k\bar H^{k,i}_t(t)$ and the drift-inference integral
$\int_0^t P^k(t)\widetilde X_t^k(r)\,\bar H_t^{k,i}(r)\,dr$.

\begin{prop}[Bounds and Lipschitz continuity]\label{prop:bounds_lip}
